\documentclass[a4paper,10pt]{article}
\usepackage[T1]{fontenc}
\usepackage[utf8]{inputenc}
\usepackage{lmodern}
\usepackage[margin=21mm]{geometry}
\usepackage{booktabs,tabularx,array,amsmath,amssymb,graphicx,xcolor,hyperref,fancyhdr}
\usepackage{enumitem}
\usepackage{float}
\definecolor{main}{HTML}{22384F}
\definecolor{muted}{HTML}{52687B}
\hypersetup{colorlinks=true,urlcolor=main,linkcolor=main}
\setlength{\parindent}{0pt}
\setlength{\parskip}{0.48em}
\setlength{\emergencystretch}{3em}
\renewcommand{\arraystretch}{1.13}
\setlist[itemize]{noitemsep,topsep=0.15em,leftmargin=1.6em}
\pagestyle{fancy}
\fancyhf{}
\fancyhead[L]{\small\color{muted}PCM: ground-truth-free PnP tuning}
\fancyhead[R]{\small\color{muted}8 November 2026}
\fancyfoot[C]{\thepage}
\setlength{\headheight}{14pt}
\begin{document}
\vspace*{0.25em}
{\LARGE\bfseries\color{main} SSDU-style parameter tuning for noisy PCM}\\[4pt]
{\large PnP-ADMM-DRUNet and GS-PnP-GS-DRUNet}\\[6pt]
{\color{muted} Experiment report \quad Version 2.0 \quad 8 November 2026}\\[2pt]
{\color{muted} Parameter search: completed and frozen \quad 450-image evaluation: complete}
\vspace{0.2cm}\hrule

\section{Executive summary}

This experiment selects plug-and-play (PnP) reconstruction parameters for noisy photocurrent mapping (PCM) \emph{without ground-truth images}. The previously reported ``oracle-tuned'' PnP baselines used ground-truth reconstruction scores on \emph{validation} images, not individual test images. We replace that supervision with a measurement-only selection loss: each candidate reconstructs from part of the acquired Hadamard measurements and predicts the withheld part. This is inspired by self-supervised learning via data undersampling (SSDU) \cite{yaman2020}, but it tunes reconstruction hyperparameters rather than neural-network weights.

We retain four distinct selections: \textbf{PnP-ADMM-tuned-5}, \textbf{PnP-ADMM-tuned-50}, \textbf{GS-PnP-tuned-5} and \textbf{GS-PnP-tuned-50}. The 5-image subset contains one validation image from each of the five PCM categories; the 50-image subset contains ten per category. All 15 ADMM and all 18 GS-PnP configurations were evaluated on both subsets, and all selections were frozen. The runner reported \textbf{no candidate failures}. The 50-image parameter search completed on tuni on \textbf{8 November 2026 at 13:42:43 UK time}.

\begin{table}[H]
\centering\small
\begin{tabularx}{\linewidth}{@{}lrlXr@{}}
\toprule
\textbf{Method} & \textbf{Images} & \textbf{Candidate} & \textbf{Chosen parameters} & \textbf{SSDU loss} \\
\midrule
PnP-ADMM & 5 & \texttt{admm\_009} & $\rho=0.5$, $\sigma=0.075$ & 0.00430443 \\
PnP-ADMM & 50 & \texttt{admm\_009} & $\rho=0.5$, $\sigma=0.075$ & 0.00451722 \\
GS-PnP & 5 & \texttt{gs\_018} & $\lambda=0.1$, $\sigma=0.075$, $\tau_0=33.3333$ & 0.00436543 \\
GS-PnP & 50 & \texttt{gs\_018} & $\lambda=0.1$, $\sigma=0.075$, $\tau_0=33.3333$ & 0.00443193 \\
\bottomrule
\end{tabularx}
\caption{Frozen results of the 5-image and 50-image tuning. SSDU loss is held-out measurement mean-squared error, \emph{not} image MSE. The losses use different validation samples.}
\label{tab:winners}
\end{table}

The two validation sizes selected \textbf{the same parameters} for each PnP algorithm, within the fixed search grids. This is evidence of selection stability, not by itself of high reconstruction quality. The \textbf{450-image test-set evaluation is complete}; Section~5 compares the frozen choices with their supervised-validation baselines using PSNR, SSIM, HaarPSI and MSE.

\section{Tuning algorithm}

\subsection*{Acquisition and four-fold splitting}

For a noisy PCM acquisition, write
\begin{equation}
 y=Ax+\varepsilon,\qquad n=256^2,\qquad m=13{,}107.
\end{equation}
Here $A$ is the \emph{normalised} dyadic Hadamard sampling operator ($J=8$) with 20\% multilevel sampling and a fully included coarse block ($J_c=6$). We use one archived NPL-default noisy acquisition per validation image; the original measurements are stored in float64 precision, while existing PnP solvers use float32 where necessary.

For each image $i$, keep all \textbf{4,096 coarse measurements} in the reconstruction set. Split only the other \textbf{9,011 acquired fine-scale measurements} into four deterministic, level-stratified held-out folds $\Lambda_{ik}$ of approximately 2,253 measurements each (blocks of 128 row positions). Let $\Theta_{ik}$ contain all acquired measurements except the held-out fold. Thus $\Theta_{ik}\cap\Lambda_{ik}=\varnothing$. The held-out measurements are not passed to the solver.

\subsection*{Prediction loss and parameter selection}

For candidate parameters $\theta$, reconstruct using the fit measurements:
\begin{equation}
 \widehat x_{\theta,ik}=\mathcal{R}_{\theta}(A_{\Theta_{ik}},y_{\Theta_{ik}}).
\end{equation}
Then calculate the loss \emph{only on the withheld measurements}:
\begin{equation}
 \ell_{ik}(\theta)=\frac{1}{|\Lambda_{ik}|}\left\|A_{\Lambda_{ik}}\widehat x_{\theta,ik}-y_{\Lambda_{ik}}\right\|_2^2.
\end{equation}
With $S_c$ the images belonging to category $c$, the selection criterion is
\begin{equation}
 \boxed{\displaystyle
 L_S(\theta)=\frac15\sum_{c=1}^{5}\frac1{|S_c|}\sum_{i\in S_c}\frac14\sum_{k=1}^{4}\ell_{ik}(\theta),\qquad
 \widehat\theta_S=\operatorname*{arg\,min}_{\theta\in\mathcal G}L_S(\theta).}
 \label{eq:ssdu}
\end{equation}
This gives every category equal weight. The $S=5$ and $S=50$ experiments use the \emph{same} candidate grid $\mathcal G$ and the same held-out folds for overlapping cases. The 5-image selections were frozen before the 50-image selections. The fixed tie tolerance was $10^{-10}$ in absolute loss, with grid order breaking ties.

\subsection*{Candidate grids and solver details}

\begin{itemize}
\item \textbf{PnP-ADMM}: $\rho\in\{0.1,0.25,0.5,1,2\}$ and denoiser $\sigma\in\{0.03,0.05,0.075\}$, giving \textbf{15 candidates}. The solver permits 20--100 iterations, absolute residual tolerance $10^{-4}$, relative residual tolerance $10^{-3}$, and requires five consecutive passes to stop.
\item \textbf{GS-PnP}: $\lambda\in\{0.01,0.03,0.1\}$, $\sigma\in\{0.03,0.05,0.075\}$, and independent initial step size $\tau_0\in\{10,100/3\}$, giving \textbf{18 candidates}. It allows 100 accepted iterations, 20 backtracking attempts, $\gamma=0.1$, $\eta=0.9$ and stopping tolerance $10^{-6}$. The reconstruction output is the \emph{paper-final $z$} image, not the internal $x$ iterate \cite{hurault2022}.
\end{itemize}

The fidelity is $\frac12\|A_{\Theta}x-y_{\Theta}\|_2^2$, with no additional division by the fit-set size. The operator audit confirmed the full and split Hadamard projections agree. The GS-PnP implementation initially showed GPU-dependent backtracking differences; after deterministic CUDA settings were enabled, the new adapter matched the frozen earlier GS solver exactly in the equivalence test.

\subsection*{Limits of the measurement-only criterion}

The simulated PCM noise can be correlated across acquisition times. If the fit and held-out noise are correlated, the held-out prediction error is not necessarily an unbiased estimate of clean measurement prediction risk. Also, if $h\in\ker(A)$, then $A(x+h)=Ax$: unobserved fine detail cannot be assessed from acquired measurements alone. Finally, an SSDU fit fold contains fewer rows than the final 20\% acquisition. The experiment keeps the fidelity convention fixed and transfers the selected parameter values directly to the full acquisition. These limitations must be considered when later comparing image metrics.

\section{Step-by-step plan}

\small
\begin{tabularx}{\linewidth}{@{}p{8mm}X p{22mm}@{}}
\toprule
\textbf{Step} & \textbf{Action} & \textbf{Status} \\
\midrule
0 & Verify the archived measurements; freeze five-image subset, 50-image set, candidate grid and 200 measurement-fold files. & Completed \\
1 & Check normalisation, fit/holdout disjointness, numerical consistency and deterministic solver adapters. & Completed \\
2 & Score all 15 ADMM and 18 GS-PnP candidates on five images, using four folds each. & Completed \\
3 & Independently freeze PnP-ADMM-tuned-5 and GS-PnP-tuned-5, including all candidate scores. & Completed \\
4 & Evaluate the identical candidate grids using all 50 validation images; reuse saved five-image folds. & Completed \\
5 & Verify all 50-image candidates and freeze the two tuned-50 configurations. & Completed \\
6 & Run all four SSDU methods and seven existing baselines on the fixed 450-image test acquisitions. & Not yet run \\
7 & Produce combined PSNR, SSIM, HaarPSI, MSE, timing and paired-comparison figures. & Not yet run \\
8 & Update this document with image quality results, failure analysis and reproducibility notes. & Future update \\
\bottomrule
\end{tabularx}
\normalsize

\section{Experiment running}

\subsection*{Data and reproducibility}

The validation input is the \textbf{frozen acquisition-only archive} from the earlier oracle WLS experiment. It contains exactly \textbf{50 validation acquisitions and no test acquisitions}: 10 from each of CIGS, Organic, Perovskite, Reference and Silicone, under 20\% sampling with NPL-default noise. The pilot image IDs are CIGS/0733, Organic/0008, Perovskite/0283, Reference/0768 and Silicone/0589 (validation permutation positions 901, 903, 905, 907 and 909). The source archive manifest begins with SHA256 \texttt{b30ff71a305bead9}.

Tuning ran on \textbf{tuni}. PnP-ADMM used the RTX 4070 and GS-PnP used the RTX 4090. Deterministic cuDNN and cuBLAS settings were enabled for numerical repeatability. The first five cases per grid were reused in full50 instead of recomputed. Across the 33 parameter configurations and 50 images with four folds each, the planned count was \textbf{6,600 unique fold reconstructions}. Each fold reconstruction can take many ADMM or GS iterations, and GS backtracking can require additional network evaluations.

The search was saved incrementally to JSON. Each candidate has per-image, per-fold losses, stopping diagnostics, aggregate category means and checksums. An unattended supervisor resumed the remaining full50 candidates after the initial single-candidate run, checked for incomplete candidates and failures, then generated frozen winner records. It finished at \textbf{13:42:43 UK time on 8 November 2026}. The numerical code revision recorded at run start was \texttt{c8842f5c42e8}; the grids, original acquisition manifest and code source files have hashes stored in the experiment archive.

\subsection*{Audit results}

The fit and held-out projections matched their corresponding rows of the full operator in the numerical audit. Adjoint relative errors were 0 for the fit operator and about $9.4\times10^{-8}$ for the held-out operator. Under deterministic settings, the GS adapter matched the previously frozen GS reconstruction exactly (maximum absolute difference 0). The ADMM adapter completed its iteration-bound smoke test. All \textbf{33 candidates} were completed in both the five-image and 50-image searches, with \textbf{no recorded candidate failures}. No ground-truth image or test-set acquisition was read for candidate selection.

\clearpage
\section{Results and deliverables}

\subsection*{Candidate ranking}

\begin{figure}[H]
\centering
\includegraphics[width=.77\linewidth]{figures/candidate_losses.pdf}
\caption{SSDU measurement-holdout losses over all candidate configurations. The two series use five and 50 validation images, respectively. Lower loss indicates better held-out measurement prediction, not necessarily higher PSNR.}
\end{figure}

\begin{table}[H]
\centering\footnotesize
\begin{tabularx}{\linewidth}{@{}l c X@{}}
\toprule
\textbf{Algorithm} & \textbf{Images} & \textbf{Three best candidates by SSDU loss} \\
\midrule
ADMM & 5 & \texttt{admm\_009} (0.00430443), \texttt{admm\_006} (0.00436202), \texttt{admm\_012} (0.00467883) \\
ADMM & 50 & \texttt{admm\_009} (0.00451722), \texttt{admm\_006} (0.00463695), \texttt{admm\_012} (0.00492435) \\
GS-PnP & 5 & \texttt{gs\_018} (0.00436543), \texttt{gs\_016} (0.00437685), \texttt{gs\_012} (0.00439187) \\
GS-PnP & 50 & \texttt{gs\_018} (0.00443193), \texttt{gs\_012} (0.00446876), \texttt{gs\_017} (0.00447703) \\
\bottomrule
\end{tabularx}
\caption{The primary, ground-truth-free ranking uses equal-category means of four-fold held-out measurement MSE.}
\end{table}

\subsection*{Per-category measurement error}

\begin{table}[H]
\centering\footnotesize
\begin{tabular}{@{}lrrrr@{}}
\toprule
& \multicolumn{2}{c}{\textbf{ADMM winner}} & \multicolumn{2}{c}{\textbf{GS-PnP winner}} \\
\cmidrule(lr){2-3}\cmidrule(lr){4-5}
\textbf{Category} & \textbf{5 images} & \textbf{50 images} & \textbf{5 images} & \textbf{50 images} \\
\midrule
CIGS & 0.003846 & 0.003747 & 0.003929 & 0.003773 \\
Organic & 0.004274 & 0.004375 & 0.004273 & 0.004319 \\
Perovskite & 0.004348 & 0.004085 & 0.004697 & 0.004135 \\
Reference & 0.004097 & 0.004871 & 0.004068 & 0.004770 \\
Silicone & 0.004957 & 0.005508 & 0.004860 & 0.005163 \\
\bottomrule
\end{tabular}
\caption{Category-mean held-out measurement MSE for the selected configurations. The five-image stage has one image per category; the 50-image stage has ten.}
\end{table}

\begin{figure}[H]
\centering
\includegraphics[width=.82\linewidth]{figures/category_losses.pdf}
\caption{Variation of held-out prediction error across the five device categories, for the frozen winner at each validation-set size.}
\end{figure}

\subsection*{Deliverable files on tuni}

The editable report \texttt{SSDU\_PnP\_Tuning\_Report.tex}, the compiled \texttt{SSDU\_PnP\_Tuning\_Report.pdf} and the two vector figures are kept together under the SSDU experiment's \texttt{report/} directory. The report can be edited and recompiled there after the 450-image evaluation has been incorporated.

The numerical evidence remains in the same experiment tree: \texttt{prepared/} contains acquisition and fold hashes, the grid and numerical audits; \texttt{pilot5/} and \texttt{full50/} contain all candidate scores and frozen parameter selections; \texttt{full50/runner/} contains execution logs. The Git branch \texttt{ssdu-pnp-tuning-20261108} stores the measurement-only scoring and reconstruction code. The intended 11-method test comparison is recorded in \texttt{comparison\_methods.json}. No source images or archived acquisitions were modified.


\subsection*{Matched 450-image test: all eleven methods}
We reconstructed the fixed 20\% NPL-noisy acquisitions of all 450 test images, 90 per category.
The seven existing comparison methods use verified historical numerical scores for
zero-filled, SPGL1-BP, SPGL1-BPDN, previously oracle-tuned ADMM, scalar TV
($\lambda=0.011$, 1024 iterations) and original UTV-C60, with the same
measurement/mask SHA256 hashes as the new runs. The prior Step-6 GS-PnP
oracle was newly reconstructed as \emph{paper-final $z$}, not substituted by
the separate historical GS-v2 method. All four SSDU selections were frozen
before reading these test metrics, and no test-based parameter adjustment was made.

The two SSDU ADMM selections coincide exactly with the earlier validation-tuned oracle settings $\rho=0.5$, $\sigma=0.075$ and its stopping rule. Therefore their archived 450 test scores are \emph{identical by construction}, not independent repeated runs.

The two SSDU GS-PnP selections also coincide, so they share one actual 450-image reconstruction run. The separate prior frozen Step-6 GS-PnP oracle was reconstructed on the same matched acquisition set.

\begin{table}[htbp]
\centering\scriptsize
\begin{tabular}{@{}lrrrr@{}}
\toprule
\textbf{Method} & \textbf{PSNR, dB} & \textbf{SSIM} & \textbf{HaarPSI} & \textbf{MSE}\\
\midrule
Zero-filled & 24.97 $\pm$ 1.97 & 0.4823 $\pm$ 0.0527 & 0.5513 $\pm$ 0.0520 & 3.526e-03 $\pm$ 1.623e-03 \\
SPGL1-BP & 25.42 $\pm$ 1.53 & 0.4353 $\pm$ 0.0354 & 0.6452 $\pm$ 0.0772 & 3.061e-03 $\pm$ 1.184e-03 \\
SPGL1-BPDN & 29.37 $\pm$ 3.41 & 0.7161 $\pm$ 0.1036 & 0.6581 $\pm$ 0.0994 & 1.552e-03 $\pm$ 1.185e-03 \\
PnP-ADMM-oracle & 35.75 $\pm$ 3.43 & 0.8970 $\pm$ 0.0770 & 0.8814 $\pm$ 0.0619 & 3.536e-04 $\pm$ 2.433e-04 \\
GS-PnP-oracle & 35.11 $\pm$ 3.64 & 0.8709 $\pm$ 0.0871 & 0.8628 $\pm$ 0.0735 & 4.184e-04 $\pm$ 2.866e-04 \\
TV (val.-tuned) & 30.03 $\pm$ 1.94 & 0.7212 $\pm$ 0.0553 & 0.7844 $\pm$ 0.0504 & 1.103e-03 $\pm$ 5.472e-04 \\
UTV-C60 & 32.30 $\pm$ 2.93 & 0.8450 $\pm$ 0.0882 & 0.7699 $\pm$ 0.1101 & 7.279e-04 $\pm$ 4.532e-04 \\
PnP-ADMM-tuned-5 & 35.75 $\pm$ 3.43 & 0.8970 $\pm$ 0.0770 & 0.8814 $\pm$ 0.0619 & 3.536e-04 $\pm$ 2.433e-04 \\
PnP-ADMM-tuned-50 & 35.75 $\pm$ 3.43 & 0.8970 $\pm$ 0.0770 & 0.8814 $\pm$ 0.0619 & 3.536e-04 $\pm$ 2.433e-04 \\
GS-PnP-tuned-5 & 32.59 $\pm$ 2.70 & 0.7570 $\pm$ 0.0967 & 0.8208 $\pm$ 0.0775 & 6.628e-04 $\pm$ 3.910e-04 \\
GS-PnP-tuned-50 & 32.59 $\pm$ 2.70 & 0.7570 $\pm$ 0.0967 & 0.8208 $\pm$ 0.0775 & 6.628e-04 $\pm$ 3.910e-04 \\
\bottomrule
\end{tabular}
\caption{Mean $\pm$ sample standard deviation of 450 paired reconstructions.
Scores computed on raw (unclipped) image arrays. Identical parameter sets share
the same test computation but retain all labels.}
\end{table}

\begin{figure}[htbp]
\centering
\includegraphics[width=\linewidth]{../test450/v1/figures/all_450_11_method_boxplots.pdf}
\caption{All 11 methods shown together over the 450 fixed test acquisitions.
The box plots exclude plotted outliers for legibility but include every case
in the descriptive statistics.}
\end{figure}

\begin{figure}[htbp]
\centering
\includegraphics[width=\linewidth]{../test450/v1/figures/all_450_11_method_mean_sd.pdf}
\caption{All 11 methods: test-set average and sample standard deviation for
PSNR, SSIM, HaarPSI and MSE.}
\end{figure}

The mean GS-PnP SSDU minus supervised-validation GS-PnP oracle PSNR difference is -2.519 dB (sample SD 1.890 dB). For ADMM the difference is exactly zero because the parameter and numerical solver choices coincide.

\subsection*{Reconstruction time}
The table gives seconds per reconstructed image, from new GPU runs or verified
historical records. Offline tuning time is excluded. The methods were not
all timed on the same GPU; therefore these values are descriptive, rather than
a controlled hardware-speed benchmark. Identical 5/50 selections reuse one
numerical reconstruction, not two independent computations.

\begin{table}[htbp]
\centering\small
\begin{tabular}{@{}lrl@{}}
\toprule
\textbf{Method} & \textbf{Timed cases} & \textbf{Mean $\pm$ SD, seconds/image}\\
\midrule
Zero-filled & 0 & not archived \\
SPGL1-BP & 450 & 1.025 $\pm$ 0.025 \\
SPGL1-BPDN & 450 & 0.249 $\pm$ 0.089 \\
PnP-ADMM-oracle & 450 & 0.421 $\pm$ 0.176 \\
GS-PnP-oracle & 450 & 1.329 $\pm$ 0.483 \\
TV (val.-tuned) & 450 & 0.683 $\pm$ 0.005 \\
UTV-C60 & 0 & not archived \\
PnP-ADMM-tuned-5 & 450 & 0.421 $\pm$ 0.176 \\
PnP-ADMM-tuned-50 & 450 & 0.421 $\pm$ 0.176 \\
GS-PnP-tuned-5 & 450 & 2.016 $\pm$ 1.099 \\
GS-PnP-tuned-50 & 450 & 2.016 $\pm$ 1.099 \\
\bottomrule
\end{tabular}
\caption{Available reconstruction times. Missing timings are not estimated.}
\end{table}

\subsection*{Per-category comparison}
Table~\ref{tab:categorypsnr} summarises PSNR means for the seven numerical
choices that are distinct after accounting for the repeated SSDU selections.
The full 11-method per-category metrics, standard deviations and boxplots
are retained as CSV, JSON and a five-page PDF in the test450 results directory.

\begin{table}[htbp]
\centering\tiny
\resizebox{\linewidth}{!}{%
\begin{tabular}{@{}lrrrrrrrr@{}}
\toprule
\textbf{Category} & \textbf{ZF} & \textbf{SPGL1} & \textbf{BPDN} &
\textbf{ADMM} & \textbf{GS oracle} & \textbf{TV} &
\textbf{C60} & \textbf{GS SSDU}\\
\midrule
CIGS & 25.71 & 26.49 & 31.10 & 39.30 & 37.89 & 31.16 & 33.04 & 35.31 \\
Organic & 25.49 & 26.37 & 32.05 & 38.90 & 38.94 & 31.50 & 35.72 & 34.96 \\
Perovskite & 27.06 & 26.43 & 31.60 & 35.48 & 35.35 & 31.33 & 32.61 & 31.70 \\
Reference & 24.25 & 24.39 & 26.72 & 32.14 & 31.72 & 28.30 & 29.63 & 29.73 \\
Silicone & 22.34 & 23.44 & 25.38 & 32.95 & 31.65 & 27.87 & 30.49 & 31.25 \\
\bottomrule
\end{tabular}}
\caption{Per-category mean PSNR (dB), with 90 test images in each category.
ADMM and GS SSDU repeated five/50 selections have been collapsed only in
this compact table, not in the full 11-method results.}
\label{tab:categorypsnr}
\end{table}

\section{Conclusion}

The planned \textbf{SSDU-style ground-truth-free parameter search is complete}. The five-image and 50-image parameter selections coincide for both PnP algorithms within the tested grids. The PnP-ADMM selection is $\rho=0.5$, $\sigma=0.075$; GS-PnP selects $\lambda=0.1$, $\sigma=0.075$, $\tau_0=100/3$. Both 5-image and 50-image selection records were frozen separately before the 450-image test evaluation.

This agreement suggests that five images were sufficient to select the same grid point \emph{in this experiment}. It does not establish the equivalence of SSDU tuning and supervised validation tuning for image quality. Temporally correlated PCM noise can bias measurement prediction errors, while thin features in unsampled directions cannot be assessed from the acquired rows alone.

We evaluated the four frozen SSDU configurations alongside seven pre-existing methods over 450 fixed 20\% NPL-default PCM acquisitions. Section~5 reports PSNR, SSIM, HaarPSI, MSE, runtime and per-category scores on the same noisy acquisitions. We made no parameter changes after inspecting test metrics. The wider project had previously examined these 450 test images, so this is a \emph{retrospective} controlled comparison rather than a wholly untouched prospective evaluation.


\paragraph{Test-set conclusion.} The SSDU-5 and SSDU-50 choices match within each
PnP family, so collecting ten validation images per category did not change
the selected grid point. The ADMM SSDU choice is exactly the original
ground-truth-validation ADMM oracle parameter pair; the GS-PnP SSDU
selection is distinct from the earlier Step-6 GS oracle. The mean GS-PnP SSDU minus supervised-validation GS-PnP oracle PSNR difference is -2.519 dB (sample SD 1.890 dB). For ADMM the difference is exactly zero because the parameter and numerical solver choices coincide.
The interpretation is limited by possible temporal PCM noise correlation,
the null-space of undersampled Hadamard imaging and prior exploratory use of
these 450 images in the broader project. No further test-driven tuning is allowed.

\renewcommand{\refname}{Primary references}
\begin{thebibliography}{9}
\bibitem{yaman2020} B. Yaman et al., ``Self-Supervised Learning of Physics-Guided Reconstruction Neural Networks without Fully-Sampled Reference Data'', \emph{Magnetic Resonance in Medicine}, 84:3172--3191, 2020. See the method for partitioning acquired measurements into data-consistency and loss sets. \url{https://doi.org/10.1002/mrm.28378}.
\bibitem{hurault2022} S. Hurault, A. Leclaire and N. Papadakis, ``Gradient Step Denoiser for Convergent Plug-and-Play'', \emph{ICLR}, 2022. See Algorithm 1 and the backtracking procedure. \url{https://openreview.net/pdf?id=fPhKeld3Okz}.
\end{thebibliography}
\end{document}
